\section{Limitations}
\label{sec:limitations}

\paragraph{Four datasets is the binding constraint.} The statistical unit is one outer partition, but the unit of
\emph{generalization} is the dataset, and there are four. Every replication-level interval here describes partition
sensitivity within a dataset, not uncertainty about a population of datasets, so every cross-dataset mechanism
statement is an association over four points and not a tested hypothesis. The ordering ``real anchors matter most
where the surrogate is unreliable'' is not one we can claim at all: the median anchor effect is largest on Santander,
whose ROC-AUC surfaces never pass the gate, and second largest on UCI credit, whose surfaces always pass
(\cref{sec:results:anchors}). All four datasets are binary tabular problems of $3\times10^{4}$ to $2\times10^{5}$
rows, and all four are financial-risk problems --- retail banking, insurance claims, motor insurance and consumer
credit --- three of them from Kaggle competitions, so the four points are not four independent domains. Nothing here
establishes behaviour on multiclass problems, text or image data, on very small samples, or on ensembles of more than
five members, where the simplex and design grow quickly.

\paragraph{The A-to-B contrast changes four things at once.} Replacing the surrogate anchors with the real
single-objective optima simultaneously moves the utopia point, rescales both normalization axes through the payoff
matrix, reorients the CHIM and all 66 quasi-normal directions, and --- because the three vertex values of $\beta$
return their anchor directly --- substitutes three returned candidates that are themselves points of the empirical
reference (\cref{sec:method:nbi}). It is a contrast between two \emph{specifications} of the anchor source, not a
single-factor experiment. The injection control of \cref{sec:results:anchors} separates the fourth of these from the
other three by rescoring NBI-A's own candidate set augmented with the same three anchors, and shows that the split is
dataset-dependent: the injected points alone account for a median $80\%$ of the hypervolume gap on Santander and
$90\%$ on BNP Paribas, but $49\%$ on Porto Seguro and $1\%$ on UCI credit. It does not separate the first three from
one another, and no control in this study does; doing so would require an NBI implementation in which the utopia
point, the normalization and the direction set can be varied independently.

\paragraph{The model zoo is fixed and shapes the coefficient findings.} Base learners and hyperparameters were fixed
in advance with no search: the study measures the weighting layer with the model layer held constant. The zoo
deliberately contains a poorly calibrated member and an expensive one, which makes the cost objective non-trivial and
also drives the $\beta_{ij}$ result of \cref{sec:results:coefficients}: the largest interactions attach to the weakest
vertex. A more homogeneous zoo would produce smaller interaction coefficients and a smaller vertex gap to explain, so
that finding is conditional on component heterogeneity; these data cannot say how it behaves for a zoo of near-equal
members.

\paragraph{The Pareto reference is empirical.} With no analytical front, all convergence and coverage indicators are
computed against a sample-based reference of about $10^{5}$ points per replication. Residual displacement is small but
non-zero: 116 of 120 replications met the 5\% tolerance and four stopped at the three-round cap with at most 6\%,
bounding the objective-space effect at 0.2\% of hypervolume. More importantly, the final reference is the
non-dominated union of the sampled core with all candidate sets, so methods contributing many points --- NBI-B and
NBI-C contribute 30--40\% --- are partly graded against their own output. This inflates their absolute indicator
values without affecting the paired contrasts on which every primary claim rests; absolute hypervolume ratios should
be read with that in mind.

\paragraph{The three NBI arms are not compute-matched, and NBI-C is not matched to anything.} Neither NBI-A nor NBI-B
spends real objective evaluations inside the optimizer, but NBI-B is not free: its anchors are produced by the
single-objective reference stage (\cref{sec:method:refs}), whose derivative-free ROC-AUC search alone evaluates
$4\times10^{4}$ compositions per replication. NBI-C consumes about $4.2\times10^{5}$ per replication and 2.3 to 6.2 times
the wall-clock time of NBI-B; the two cheap comparators consume about 66 evaluations and a few seconds. The comparison
is therefore between \emph{methods as specified}, not at equal budget, and supports no efficiency claim. The random
Dirichlet comparator is matched only in candidate count and is degenerate under the weighted cost; it is a floor check
on the reference geometry. The budget-matched question is answered only at the high end, by the NSGA-II baseline of
\cref{sec:results:nsga2}; the low-budget question --- what an evolutionary method or a random-search ladder achieves
on the 66 design points and 100 gate points the surrogate arms use --- is the one a surrogate paper most needs and the
one we do not answer.

\paragraph{Solver semantics differ between the arms.} NBI-A and NBI-B optimize smooth polynomials, and their
subproblems count as successful only when SLSQP certifies optimality. NBI-C optimizes a piecewise-constant ROC-AUC,
where such certification is generically unattainable, so a subproblem is accepted when its equality residual falls
below $10^{-3}$ under a reduced iteration budget. The reported ``success rates'' are therefore not on a common scale
across arms, and the UCI credit gap between NBI-C and NBI-B is substantially a convergence artifact of NBI-B rather
than evidence about surrogate quality (\cref{sec:results:metamodelfree}); we therefore report both the filtered and
unfiltered versions of that comparison.

\paragraph{The support cost is a threshold approximation.} Deployment cost is modelled by \cref{eq:supportcost} with
$\varepsilon = 10^{-3}$, fixed before the study and never changed. It is a proxy that ignores memory, batching, model
loading and shared preprocessing, and treats a weight of $1.1\times10^{-3}$ as a full model and one of
$0.9\times10^{-3}$ as free. \Cref{sec:results:conflict} shows this is not hypothetical: an AUC-inert
$k$-nearest-neighbour weight just above the threshold produced a spurious $250$~ms difference. Any support cost
therefore needs a weight-cleaning step, and conclusions that depend on it are sensitive to $\varepsilon$ and to the
cost range that sets the normalization box.

\paragraph{Cost measurements are single-machine and single-implementation.} All $c_i$ are wall-clock medians on one
workstation with a fixed library stack and batch size; they would change with hardware, batching, a compiled inference
runtime, or an approximate-nearest-neighbour index replacing brute-force search, which would remove the study's
largest cost asymmetry. The cost objective is realistic for the setting measured, not universal.

\paragraph{No claim of universal dominance, and confounded comparators.} NBI-C is best or tied-best on the primary
endpoints under the weighted cost on these four datasets; it ties with random scalarization on UCI credit, loses to it
on Santander $\igdp$, and loses its support-cost advantage on BNP Paribas and UCI credit. Its advantage over
scalarization on BNP Paribas and Porto Seguro is confounded with the objective source: NBI-C optimizes exact
objectives, scalarization optimizes surrogates. Separating the NBI construction from the objective source would
require a scalarization arm on the cached out-of-fold objectives, which we did not run.

\paragraph{Selection rules and the holdout analysis.} Holdout transfer is reported for a knee rule under the weighted
cost, with TOPSIS and the support cost as sensitivity analyses, and agreement rates differ noticeably between these
configurations. They characterize how close the candidates are relative to holdout noise, not the general reliability
of out-of-fold selection. What we verify is transfer \emph{within} a candidate set and in level. The cross-method
question --- whether the arm with the better front also yields the better selected solution --- is answered only
descriptively, and on the two contrasts that carry our Holm-significant primary results it is answered in the negative
at the selected point (\cref{sec:disc:transfer}); no statement here should be read as ``better fronts give better
picks''.

\paragraph{Only one external optimizer was evaluated.} We added a single canonical evolutionary baseline, NSGA-II at a
matched real-objective evaluation budget, and it outperformed the metamodel-free arm on every dataset. That is one
algorithm at one configuration on one problem class: the study is not a competition among multiobjective optimizers,
and nothing here ranks NSGA-II against MOEA/D, NSGA-III or a surrogate-assisted evolutionary method. Nor is the
comparison time-matched --- at equal evaluation count NSGA-II costs about $2.2\times$ the wall clock in our
implementation, and a differently batched one would shift that ratio.
